\chapter{INTRODUCTION}\label{chap:introduction}

This chapter presents the foundation of the study by establishing the
residential electrical-fire problem, the limitations of conventional
protection in detecting progressive fault conditions, and the need for a
complementary socket-level monitoring device. It also defines the
problem addressed by the research, identifies the significance of the
proposed system, states the objectives that guided the development and
evaluation of the prototype, and clarifies the scope and limitations of
the study.

\section{Background of the Study}\label{background-of-the-study}

Electrical faults remained one of the leading causes of residential
fires in the Philippines. Bureau of Fire Protection records showed that
electrical faults accounted for 27.09\% of reported fire incidents,
while related literature also identified undetected electrical
abnormalities, faulty outlets, overloaded circuits, degraded wiring, and
faulty appliances as recurring causes of residential fire events
{[}1{]}--{[}3{]}. These incidents did not only result in property loss,
but also caused injuries, displacement, and long-term economic hardship
for affected households.

Residential electrical service in the Philippines, particularly in
provincial and urban domestic settings such as those commonly found in
Naga City, generally consisted of a 230V, single-phase, two-wire, 60Hz
supply {[}4{]}. Although this configuration could technically
accommodate more advanced protective devices, the actual residential
protection landscape still depended primarily on miniature circuit
breakers (MCBs), with wider adoption of specialized protective devices
such as AFCIs and GFCIs remaining limited by cost, availability, and
lack of mandatory implementation in ordinary homes {[}5{]}. As a result,
many households still relied on protection systems that were effective
mainly against high overcurrent and short-circuit events, but not
against slowly developing pre-fire electrical conditions.

This created a critical safety gap. Conventional MCBs were designed
primarily to protect fixed building wiring from sustained overcurrent
and short-circuit damage. They were not designed to warn occupants about
gradual or localized electrical hazards before those hazards escalated
into ignition. In particular, they did not directly detect progressive
contact heating at a socket, persistent sub-trip overload behavior, or
low-current series arc faults that occurred below the breaker's
immediate trip region. In poor or degraded electrical connections, a
small region of elevated resistance could generate intense localized
heating under otherwise ordinary current flow. This condition could
persist long enough to char surrounding material, deform outlet
components, and eventually cause fire even without a classical breaker
trip {[}6{]}, {[}7{]}.

At the same time, recent developments in non-intrusive electrical
monitoring, embedded machine learning, and low-cost microcontrollers
made it possible to design a more accessible protective device.
Non-intrusive sensing enabled waveform-based monitoring without
modifying existing building wiring {[}8{]}. Machine learning methods,
especially Random Forest-based classifiers, had already demonstrated
strong performance in electrical-fault classification tasks {[}9{]}.
TinyML deployment on microcontrollers further showed that real-time
inference could be achieved within tight memory, latency, and power
constraints {[}10{]}. These developments suggested that a practical,
socket-level protective device for residential environments was already
technically feasible.

This study therefore developed an ESP32-based protective smart plug that
combined voltage, current, and thermal sensing with hybrid fault
detection. The system was designed as a complementary protection layer
for Philippine single-phase residential circuits. It focused on three
progressive fire-inducing fault conditions: overload-related hazardous
overcurrent behavior, series arc faults from unstable electrical
conduction, and contact heating at the plug-socket interface. Rather
than replacing mandatory protection, the device was intended to provide
earlier warning, localized monitoring, and automatic disconnection when
hazardous conditions were detected before they escalated into more
serious fire events.

\section{Statement of the Problem}\label{statement-of-the-problem}

The preceding discussion shows that residential electrical fire
prevention does not only depend on clearing severe short circuits or
major overloads. It also requires attention to progressive electrical
hazards that may remain energized while still producing thermal or
waveform abnormalities. This study was therefore concerned with the
protection gap between normal household operation and the point at which
conventional circuit protection is forced to operate.

Residential electrical fires remained a persistent public safety concern
because the electrical protection normally present in ordinary homes was
not sufficient to address all dangerous fault mechanisms. The Philippine
residential setup commonly relied on MCBs as the primary mandatory
protective device {[}4{]}, {[}5{]}. While effective against major
overload and short-circuit events, these devices did not provide early
warning for progressive hazards and did not directly detect the
low-current or localized fault mechanisms that often preceded electrical
fires.

A particular problem involved fault conditions that developed below the
level of obvious breaker action. A household branch circuit could remain
energized while a socket or terminal connection slowly degraded, while
intermittent series arcing occurred during unstable conduction, or while
a connected load remained in a hazardous but not immediately
breaker-tripping overcurrent state. In such cases, occupants could
remain unaware of the developing danger until visible damage, smoke, or
ignition had already begun. This protection gap became more serious in
settings where advanced commercial arc-fault protection was uncommon and
where households continued to rely on low-cost, conventional protective
hardware alone.

Although advanced protective technologies already existed, their direct
adoption in ordinary Philippine homes remained limited. AFCIs and
related commercial systems could improve safety, but they were not yet
common in many residential settings because of cost, product
availability, installation practice, and regulatory limitations {[}5{]}.
Existing research prototypes also often relied on intrusive
instrumentation, laboratory-grade measurement equipment, or
computation-heavy processing that was not suitable for low-cost embedded
deployment. Moreover, arc-fault detection systems in both literature and
practice still faced nuisance-alarm challenges when normal appliance
behavior, switching events, or electrically noisy loads produced
signatures similar to true faults.

This study therefore addressed the need for an affordable, socket-level,
non-intrusive protective device that could operate alongside
conventional circuit protection. Specifically, the device needed to
detect and respond to three progressive fault classes relevant to
residential fire safety: overload-related hazardous current behavior,
induced series arc faults, and contact heating at the socket interface.
The device also needed to operate on resource-constrained embedded
hardware, provide local protection even without cloud access, and
support logging and remote visibility when network connectivity was
available.

\section{Significance of the Study}\label{significance-of-the-study}

This study is expected to provide substantial benefits to residential
users, fire-safety stakeholders, engineering researchers, and
policymakers by addressing a practical gap in household electrical
protection. Since progressive electrical hazards may develop before
conventional protective devices operate, the development of a
socket-level smart plug offers a complementary approach for earlier
detection, localized monitoring, and automatic response. The study also
supports the principles of Sustainable Development Goal 11 by
contributing to safer and more resilient communities through an
accessible residential safety technology {[}11{]}. The key beneficiaries
and their potential gains are outlined below.

\subsection{Residential Households}\label{residential-households}

Residential households are the primary beneficiaries of this study. The
proposed smart plug provides an added layer of protection against
electrical hazards that are not easily visible to ordinary occupants,
such as sustained overload, socket-contact heating, and series arc
faults. By monitoring voltage, current, and temperature at the outlet
level, the device can alert users and disconnect the connected load when
hazardous conditions are detected. This gives households earlier
awareness of possible electrical danger before the condition develops
into visible damage, smoke, or fire.

The study is especially relevant to ordinary residential users who rely
mainly on conventional circuit breakers for electrical protection. While
breakers remain necessary, they may not directly detect localized socket
heating or low-current arc behavior. The prototype therefore
demonstrates how a plug-through device can supplement existing household
protection without requiring permanent rewiring of the home.

\subsection{Local Government Units}\label{local-government-units}

Local government units, disaster risk reduction offices, and fire-safety
stakeholders may benefit from the study because it presents a possible
technology-based approach for improving household-level fire prevention.
Since electrical faults remain a recurring cause of residential fire
incidents, devices capable of detecting early electrical abnormalities
can support broader fire-risk reduction efforts.

The study may also help fire-safety personnel and community safety
programs explain the importance of monitoring progressive electrical
hazards, not only severe short circuits. By showing that overload,
contact heating, and series arc faults can be monitored at the socket
level, the research provides a technical basis for promoting more
proactive residential electrical-safety awareness.

\subsection{Electrical Practitioners and Safety Inspectors}\label{electrical-practitioners-and-safety-inspectors}

Electrical practitioners, technicians, and safety inspectors may benefit
from the study because it highlights fault conditions that are not
always addressed by ordinary visual inspection or conventional
overcurrent protection. The prototype demonstrates how localized sensing
can identify abnormal behavior at the plug-socket interface, where
degraded contact, poor fit, corrosion, or thermal buildup may occur.

The study may also support future diagnostic practices by showing the
value of combining current behavior, temperature behavior, and
fault-event logging. Although the device is not a substitute for
professional electrical inspection, it provides a possible tool for
observing early warning patterns that may guide maintenance,
troubleshooting, or further assessment of residential electrical
installations.

\subsection{Engineering and Academic Community}\label{engineering-and-academic-community}

The engineering and academic community may benefit from the study
through its documented integration of embedded sensing, TinyML
inference, threshold-based protection, relay actuation, and
cloud-assisted monitoring in a single residential prototype. The study
provides a development workflow that includes sensor calibration,
waveform-feature extraction, series arc data collection, model training,
embedded deployment, and controlled fault evaluation.

This contribution is relevant to future research in embedded electrical
safety, smart protection devices, TinyML, non-intrusive monitoring, and
Internet-of-Things-based residential systems. It may also serve as a
reference for researchers who aim to improve arc-fault classification,
socket-heating estimation, overload response, or real-time protection on
resource-constrained microcontrollers.

\subsection{Policymakers and Standards Stakeholders}\label{policymakers-and-standards-stakeholders}

Policymakers, standards stakeholders, and residential-safety advocates
may benefit from the study because it demonstrates the feasibility of
treating intelligent fault detection as a complement to existing
electrical protection. The prototype does not replace code-mandated
devices such as circuit breakers, but it shows how additional
outlet-level protection may help address hazards that are difficult to
detect using conventional protection alone.

The findings may contribute to future discussions on affordable
residential safety devices, smart-home electrical protection, and the
gradual modernization of household fire-prevention strategies. By
presenting a locally developed prototype for Philippine residential
conditions, the study may help inform future research, policy planning,
or pilot implementation of complementary electrical-safety technologies.

\section{Objectives}\label{objectives}

This study aimed to develop an intelligent ESP32-S3-based protective
smart plug that provides complementary, socket-level protection for
Philippine single-phase residential circuits by detecting progressive
fire-inducing electrical faults through hybrid on-device sensing,
analysis, and fault response.

Specifically, the study aimed to:

\begin{enumerate}
\def\labelenumi{\arabic{enumi}.}
\item
  Design and develop an ESP32-S3-based smart plug prototype integrating
  voltage, current, and temperature sensing, embedded processing,
  relay-based load disconnection, and local audio-visual fault
  indication for residential single-phase applications;
\item
  Calibrate and validate the voltage, current, and temperature sensing
  subsystems against appropriate reference instruments to establish
  measurement accuracy suitable for threshold-based protection and
  machine learning-based analysis;
\item
  Implement a hybrid fault-detection system that uses on-device TinyML
  classification for series arc-fault detection and deterministic
  threshold-based logic for overload and socket-heating protection;
\item
  Develop a complementary protection, alert, and monitoring mechanism
  that operates alongside existing safety devices through local
  audio-visual alerts, relay disconnection, fault-event logging, and a
  Firebase-integrated Progressive Web App for real-time visualization
  and basic remote-control functions; and
\item
  Evaluate the developed system under controlled laboratory conditions
  involving overload, induced series arc fault, and
  socket/contact-heating scenarios in terms of detection behavior,
  system response, measurement reliability, and operational feasibility.
\end{enumerate}

\section{Scope and Delimitations}\label{scope-and-delimitations}

The scope and delimitations define the boundary within which the
prototype was designed, tested, and interpreted. Since the study
involved electrical safety, embedded sensing, machine learning, and
controlled fault simulation, the boundaries are separated into the
actual coverage of the work and the methodological, technical, and
practical limitations that affected the interpretation of the findings.

\subsection{Scope of the Study}\label{scope-of-the-study}

This study focused on the design, development, calibration, and
controlled evaluation of an intelligent protective smart plug for
Philippine residential electrical systems. The prototype was intended
for 230V, single-phase, two-wire, 60Hz circuits typical of ordinary
household branch connections. It was developed as a plug-through,
socket-level device that could be inserted into an existing outlet
without requiring permanent modification of the building wiring.

The study covered the complete prototype process, including hardware
integration, firmware development, signal acquisition, sensor
calibration, feature extraction, TinyML model deployment,
threshold-based protection logic, relay actuation, local notification,
cloud logging, and remote data visualization. The final system
architecture used a ZMPT101B-based voltage sensor, an ACS37035 current
sensor digitized through an MCP3204 external SAR ADC, and a TTC05103JSY
10kΩ NTC thermistor for temperature measurements. Current analysis was
executed on the ESP32-S3 Xiao using embedded sampling, windowed feature
computation, and on-device inference for real-time protection decisions.

The developed prototype monitored three primary fault conditions
relevant to residential fire risk: overload current behavior, induced
series arc faults, and contact heating at the plug-socket interface.
Within this scope, the device was implemented as a complementary safety
function rather than a replacement for breaker protection. Overload
protection used current-based warning and trip logic, socket-heating
protection used estimated socket-temperature behavior and a conservative
40°C control threshold, and series arc-fault detection used on-device
Random Forest inference based on extracted current features and
lab-generated labeled data.

The study evaluated the prototype under controlled laboratory scenarios
involving resistive overload simulation, controlled series arc
generation, and simulated contact heating using a degraded socket
sample. The emphasis of the evaluation was on technical feasibility,
detection behavior, protection response, measurement reliability, and
prototype operability in a socket-level residential use case. It did not
seek to replace utility protection, distribution-panel protection, or
full-house fault-monitoring systems.

\subsection{Methodological Limitations}\label{methodological-limitations}

The study was conducted under controlled laboratory conditions rather
than in occupied residential homes. Although this was necessary for
safety, repeatability, and controlled comparison, laboratory conditions
did not fully reproduce the variability of actual household
environments. Factors such as wiring age, long-term humidity exposure,
appliance diversity, plug fit, environmental contamination, user
behavior, and branch-circuit variability in actual homes were not
comprehensively represented. The findings therefore demonstrated
controlled technical feasibility rather than complete real-house field
performance.

The series arc-fault dataset used for training, validation, and
detection was based on controlled fault inducement rather than naturally
occurring household arc events. The arc generator used a repeatable
electrode separation and recontact mechanism to create series arcs under
known conditions, and the collected labels were derived from induced
test events. This approach was appropriate for safe data collection and
reproducibility, but naturally occurring residential arc faults may
exhibit wider variability in duration, intermittence, load interaction,
and noise than the controlled setup reproduced. Consequently, model
generalization to all real-world arc behaviors remained limited.

The training and validation scope was also limited by the load families
included in the experiments. The device was tested using selected load
categories that were safe and practical for induced-fault
experimentation. It did not comprehensively include all possible
household load types, such as sensitive SMPS loads,
phase-angle-controlled loads, or mixed electronic loads, because some of
these devices could be damaged during induced fault testing. The
resulting model performance should therefore be interpreted within the
conditions and load classes represented in the collected dataset.

\subsection{Technical Limitations}\label{technical-limitations}

The contact-heating function was constrained by the non-intrusive
thermal architecture of the prototype. The NTC thermistor did not
directly measure the true internal contact-spot temperature inside the
socket interface. Instead, it measured a delayed external thermal
response that the firmware interpreted as an estimated
socket-temperature proxy. For that reason, the implemented 40°C
threshold should be understood as a conservative, device-specific
control threshold for an indirect external sensor that could under-read
the true hidden hotspot. This limitation was influenced by thermistor
placement, thermal lag, enclosure conduction, and the presence of
cooling hardware inside the prototype.

The prototype also remained susceptible to nuisance detections under
certain operating conditions. As with AFCI research and other
arc-detection systems, normal events such as switching transitions,
relay actuation, brushed-motor behavior, intermittent appliance
conduction, or electrically noisy loads could generate arc-like
signatures. Although the firmware and model pipeline were designed to
reduce false positives, nuisance alarms and classification ambiguity
remained possible, particularly for unseen load types or fault-adjacent
events not fully represented in the training data.

The device provided only local, socket-level sensing. It monitored the
branch segment and connected load at the outlet where it was installed;
it was not designed to provide whole-house monitoring, multi-branch
coordination, or exact upstream fault localization. It therefore could
not be treated as a substitute for panel-level protective
infrastructure, wiring inspection, or full-building electrical
diagnostics. In addition, very low-power loads were outside the reliable
operating range of the arc-detection pipeline because low current levels
may reduce signal-to-noise ratio and weaken waveform-feature
reliability.

Remote visualization, cloud logging, and web-based control also depended
on Wi-Fi and internet availability. Core protection functions were
implemented locally on the device, but convenience features such as
Firebase logging, Progressive Web App visualization, and remote
monitoring were not fully available when network connectivity was
unstable or unavailable.

\subsection{Practical Limitations}\label{practical-limitations}

The prototype was constrained by practical hardware and packaging
limitations. Because it used a relay, heatsink, fan, sensing modules,
and protective enclosure, the device was larger and heavier than an
ordinary household plug accessory and was not yet optimized for
consumer-ready ergonomics. The available GPIO, memory, enclosure space,
and power arrangement also limited future expansion of hardware
features.

The study was also limited to prototype performance and technical
feasibility. It did not comprehensively address long-term field aging,
certification, mass-manufacturing readiness, regulatory approval, market
adoption, user acceptance, maintenance practice, or full economic
comparison with commercial AFCI/GFCI solutions. Latched fault handling
also required user awareness and intervention for confirmation, reset,
or continued operation after certain protective events.

For these reasons, the prototype was not presented as a complete
solution to residential electrical fire prevention. It was instead
framed as a complementary protective layer that could operate alongside
existing technologies and provide earlier local awareness of progressive
fire-inducing electrical faults under the tested conditions.

